% s03-4-1-construction.tex —— 3.4 引言、3.4.1 与 3.4.2（印刷页 103–110）

\section{Iwahori 基本局部等价}
\label{sec:ifle}

\subsection{IFLE 的构造}
\label{sec:ifle-construction}

本小节将在 $\FactModCat_{x_0}$ 中构造 Iwahori 基本局部等价
\[
\IFLE\colon (\KL(G)_{\crit},\IKL(G)_{\crit})\longrightarrow
\bigl(\IndCoh^{*}(\Op^{\monfree}_{\check G}),
\IndCoh^{*}(\Op^{\mer}_{\check G}
\times_{\LS^{\mer}_{\check G}}\check{\mathfrak{n}}/\check B)\bigr)
\]
作为对的态射。
\srcpage{103}

回忆有作用
\[
\act\colon \Rep(\check G)\otimes \KL(G)_{\crit}\longrightarrow \KL(G)_{\crit},
\]
它来自 $\Sph_G$ 在 $\KL(G)_{\crit}$ 上的作用与幺半函子
\[
\Sat^{\mathrm{nv}}_G\colon \Rep(\check G)\longrightarrow
\Sph^{\spec}_{\check G}\simeq \Sph_G .
\]
有
\[
\act(R_{\check G,\Op}\boxtimes V^{\mathrm{ac}})\simeq V^{\mathrm{ac}},
\]
故
\[
R_{\check G,\Op}\boxtimes V^{\mathrm{ac}}\longrightarrow \coact(V^{\mathrm{ac}}).
\]
可以把作用 $\act$ 延拓为
\[
(\Rep(\check G),\QCoh(\check{\mathfrak{n}}/\check B))\otimes
(\KL(G)_{\crit},\IKL(G)_{\crit})\longrightarrow
(\KL(G)_{\crit},\IKL(G)_{\crit}),
\]
办法是把 $(\Sph_G,\Aff_G)$ 在 $(\KL(G)_{\crit},\IKL(G)_{\crit})$ 上的作用与对的幺半态射
\[
\mathrm{B}^{\mathrm{nv}}\colon
(\Rep(\check G),\QCoh(\check{\mathfrak{n}}/\check B))\longrightarrow
(\Sph^{\spec}_{\check G},\Aff^{\spec}_{\check G})\simeq (\Sph_G,\Aff_G)
\]
复合。

现在，仿照 \cite[第 5.3 节]{ABC+24a}，考虑上述作用的右伴随
\[
(\KL(G)_{\crit},\IKL(G)_{\crit})\longrightarrow
(\Rep(\check G),\QCoh(\check{\mathfrak{n}}/\check B))\otimes
(\KL(G)_{\crit},\IKL(G)_{\crit}).
\]
\srcpage{104}
它经
\begin{align*}
&(\KL(G)_{\crit},\IKL(G)_{\crit})\longrightarrow
\coact(V^{\mathrm{ac}})\text{-}\FactMod\\
&\qquad\qquad\bigl[(\Rep(\check G),\QCoh(\check{\mathfrak{n}}/\check B))\otimes
(\KL(G)_{\crit},\IKL(G)_{\crit})\bigr]
\end{align*}
分解，而沿 $R_{\check G,\Op}\boxtimes V^{\mathrm{ac}}\to \coact(V^{\mathrm{ac}})$ 的因子化
限制给出到
\[
(R_{\check G,\Op}\boxtimes V^{\mathrm{ac}})\text{-}\FactMod
\bigl[(\Rep(\check G),\QCoh(\check{\mathfrak{n}}/\check B))\otimes
(\KL(G)_{\crit},\IKL(G)_{\crit})\bigr]
\]
的映射，它等同于
\[
R_{\check G,\Op}\text{-}\FactMod(\Rep(\check G),\QCoh(\check{\mathfrak{n}}/\check B))
\otimes V^{\mathrm{ac}}\text{-}\FactMod(\KL(G)_{\crit},\IKL(G)_{\crit}),
\]
换言之，等同于
\[
R_{\check G,\Op}\text{-}\FactMod(\Rep(\check G),\QCoh(\check{\mathfrak{n}}/\check B))
\otimes (\KL(G)_{\crit},\IKL(G)_{\crit}).
\]
然而由定理~\ref{thm:3.1.5.2}，存在对的态射
\begin{align*}
&\bigl(\IndCoh^{*}(\Op^{\monfree}_{\check G}),
\IndCoh^{*}(\Op^{\mer}_{\check G}
\times_{\LS^{\mer}_{\check G}}\check{\mathfrak{n}}/\check B)\bigr)\\
&\qquad\longrightarrow
R_{\check G,\Op}\text{-}\FactMod(\Rep(\check G),\QCoh(\check{\mathfrak{n}}/\check B)),
\end{align*}
它在每个分量的最终余连通部分上是等价。

于是可以得到从
\[
(\KL(G)_{\crit},\IKL(G)_{\crit})
\]
到
\[
\bigl(\IndCoh^{*}(\Op^{\monfree}_{\check G}),
\IndCoh^{*}(\Op^{\mer}_{\check G}
\times_{\LS^{\mer}_{\check G}}\check{\mathfrak{n}}/\check B)\bigr)
\otimes (\KL(G)_{\crit},\IKL(G)_{\crit})
\]
的映射。
\srcpage{105}

由于在第一个分量上，如该文所述，
\[
\KL(G)_{\crit}\longrightarrow \Rep(\check G)\otimes \KL(G)_{\crit}
\]
是 t-正合的，故
\[
\KL(G)^{c}_{\crit}\subset \KL(G)^{b}_{\crit}
\]
被送到
\[
R_{\check G,\Op}\text{-}\FactMod(\Rep(\check G))\otimes \KL(G)_{\crit}
\]
的最终余连通部分。在第二个分量上，
\[
\Aff^{\spec}_{\check G}\otimes \IKL(G)_{\crit}\longrightarrow \IKL(G)_{\crit}
\]
对每个固定的有界对象 $\IKL(G)_{\crit}$ 是有界于右的。于是
\[
\QCoh(\check{\mathfrak{n}}/\check B)\otimes \IKL(G)_{\crit}\longrightarrow \IKL(G)_{\crit}
\]
亦然，从而右伴随
\[
\IKL(G)_{\crit}\longrightarrow \QCoh(\check{\mathfrak{n}}/\check B)\otimes \IKL(G)_{\crit}
\]
对每个固定的有界对象 $\IKL(G)_{\crit}$ 有界于左。因此
\[
\IKL(G)^{c}_{\crit}\subset \IKL(G)^{b}_{\crit}
\]
被送到
\[
R_{\check G,\Op}\text{-}\FactMod(\QCoh(\check{\mathfrak{n}}/\check B))
\otimes \IKL(G)_{\crit}
\]
的最终余连通部分。
\srcpage{106}

该态射定义了
\[
\bigl(\IndCoh^{*}(\Op^{\monfree}_{\check G}),
\IndCoh^{*}(\Op^{\mer}_{\check G}
\times_{\LS^{\mer}_{\check G}}\check{\mathfrak{n}}/\check B)\bigr)
\]
在 $(\KL(G)_{\crit},\IKL(G)_{\crit})$ 上的余作用，换言之，
\[
\bigl(\IndCoh^{!}(\Op^{\monfree}_{\check G}),
\IndCoh^{!}(\Op^{\mer}_{\check G}
\times_{\LS^{\mer}_{\check G}}\check{\mathfrak{n}}/\check B)\bigr)
\]
的作用。

与 \cite[第 6.1 节]{ABC+24a} 平行，可以考虑对的态射
\begin{equation}
(\KL(G)_{\crit},\IKL(G)_{\crit})\longrightarrow
(\Whit^{!}(\hat{\mathfrak{g}}\text{-}\mathrm{mod}_{\mathrm{crit}}),
\Whit^{!}(\hat{\mathfrak{g}}\text{-}\mathrm{mod}_{\mathrm{crit}}))
\label{eq:3.4.1}
\end{equation}
\begin{equation}
\xrightarrow{\ \mathrm{DS}^{\mathrm{enh,rfnd}}\ }
\bigl(\IndCoh^{*}(\Op^{\mer}_{\check G}),
\IndCoh^{*}(\Op^{\mer}_{\check G})\bigr),
\label{eq:3.4.2}
\end{equation}
把它看作
$\bigl(\IndCoh^{!}(\Op^{\mer}_{\check G}),
\IndCoh^{!}(\Op^{\mer}_{\check G})\bigr)$-模（在幺半意义下）的态射。于是得到对的态射，从
\[
(\KL(G)_{\crit},\IKL(G)_{\crit})
\]
到
\[
\Fun_{\IndCoh^{!}(\Op^{\mer}_{\check G})}
\Bigl(\bigl(\IndCoh^{!}(\Op^{\monfree}_{\check G}),
\IndCoh^{!}(\Op^{\mer}_{\check G}
\times_{\LS^{\mer}_{\check G}}\check{\mathfrak{n}}/\check B)\bigr),
\IndCoh^{*}(\Op^{\mer}_{\check G})\Bigr),
\]
后者等同于
\[
\bigl(\IndCoh^{*}(\Op^{\monfree}_{\check G}),
\IndCoh^{*}(\Op^{\mer}_{\check G}
\times_{\LS^{\mer}_{\check G}}\check{\mathfrak{n}}/\check B)\bigr).
\]
这个对的态射就是我们将记作 IFLE 的态射。
\srcpage{107}

\begin{remark}\label{rem:3.4.1.1}
IFLE 把 $(\Sph_G,\Aff_G)$ 在 $(\KL(G)_{\crit},\IKL(G)_{\crit})$ 上的作用与
$(\Sph^{\spec}_{\check G},\Aff^{\spec}_{\check G})$ 在
$\bigl(\IndCoh^{*}(\Op^{\monfree}_{\check G}),
\IndCoh^{*}(\Op^{\mer}_{\check G}\times_{\LS^{\mer}_{\check G}}
\check{\mathfrak{n}}/\check B)\bigr)$ 上的作用交织起来（更精确地说，带有交织数据），
交织所用的是 $\FactModCat_{x_0}$ 中代数对象的幺半等价
\[
\mathrm{B}\colon
(\Sph_G,\Aff_G)\simeq (\Sph^{\spec}_{\check G},\Aff^{\spec}_{\check G}).
\]
原因是：\cite[附录 E.10.4 至 E.11 末]{ABC+24a} 的全部内容都可以逐字搬到对的情形
（其中引理 E.10.3 换成我们对 B 的构造）。
\end{remark}

\subsection{归约到逐点等价}
\label{sec:ifle-reduction}

下面说明：要证 IFLE 是对之间的等价，只需证它在 $x_0$ 处是等价。我们的策略与
\cite[第 6.2 节]{ABC+24a} 中处理 FLE 时类似。

本小节余下部分，我们假设

\begin{proposition}\label{prop:3.4.2.1}
\[
\IFLE_{x_0}\colon \IKL(G)_{\crit,x_0}\longrightarrow
\IndCoh^{*}\bigl(\Op^{\mer}_{\check G}
\times_{\LS^{\mer}_{\check G}}\check{\mathfrak{n}}/\check B\bigr)_{x_0}
\]
\srcpage{108}
是等价。
\end{proposition}

本命题将在本节后面几个小节中证明。

注意在 $\Ran_{x_0}$ 的每个层片 $X^{(n),\circ}_{x_0}$ 上，
\[
\IFLE_{X^{(n),\circ}_{x_0}}\colon
(\IKL(G)_{\crit})_{X^{(n),\circ}_{x_0}}\longrightarrow
\IndCoh^{*}\bigl(\Op^{\mer}_{\check G}
\times_{\LS^{\mer}_{\check G}}\check{\mathfrak{n}}/\check B\bigr)
_{X^{(n),\circ}_{x_0}}
\]
可等同于
\begin{align}
\FLE_{X^{(n),\circ}}\boxtimes \IFLE_{x_0}\colon\quad
&(\KL(G)_{\crit})_{X^{(n),\circ}}\boxtimes \IKL(G)_{\crit,x_0}
\longrightarrow\label{eq:3.4.3}\\
&\IndCoh^{*}(\Op^{\monfree}_{\check G})_{X^{(n),\circ}}
\boxtimes \IndCoh^{*}\bigl(\Op^{\mer}_{\check G}
\times_{\LS^{\mer}_{\check G}}\check{\mathfrak{n}}/\check B\bigr)_{x_0},
\label{eq:3.4.4}
\end{align}
它是等价，因为 FLE 是等价。

于是只需说明 $\IKL(G)_{\crit}$ 紧生成、且 IFLE 保持紧对象。这将表明 IFLE 作为范畴层
之间的态射有连续右伴随，从而是否为等价可以逐层片地检验。

正如 $\KL(G)_{\crit}$ 由 $\Rep(L^+G)$ 的紧生成元在
\[
\mathrm{ind}(\hat{\mathfrak{g}}^{\mathrm{crit}},L^+G)_{L^+G}\colon
\Rep(L^+G)\longrightarrow \KL(G)_{\crit}
\]
下的像生成，范畴 $\IKL(G)_{\crit}$ 由因子化 $\Rep(L^+G)$-模范畴 $\Rep(I)$ 的紧生成元在
\[
\mathrm{ind}(\hat{\mathfrak{g}}^{\mathrm{crit}},I)_{I}\colon
\Rep(I)\longrightarrow \IKL(G)_{\crit}
\]
下的像生成；后者可以通过模仿 \cite[附录 B.14]{ABC+24a} 的构造来验证。

\begin{remark}\label{rem:3.4.2.2}
下面的论证由 Lin Chen 向我们解释。
\end{remark}

\begin{lemma}\label{lem:3.4.2.3}
对象 $c\in \IndCoh^{*}\bigl(\Op^{\mer}_{\check G}
\times_{\LS^{\mer}_{\check G}}\check{\mathfrak{n}}/\check B\bigr)$
是紧的，只要对每个 $\lambda\in \Lambda_G$，
\srcpage{109}
沿
\[
\Op^{\mer}_{\check G}\times_{\LS^{\mer}_{\check G}}\check{\mathfrak{n}}/\check B
\longrightarrow \Op^{\mer}_{\check G}
\]
对
\[
c(\lambda)=c\otimes
\bigl(\Op^{\mer}_{\check G}\times_{\LS^{\mer}_{\check G}}
\check{\mathfrak{n}}/\check B\to \check{\mathfrak{n}}/\check B\bigr)^{!}\mathcal{O}(\lambda)
\]
的 $*$-推前在 $\IndCoh^{*}(\Op^{\mer}_{\check G})$ 中是紧的。
\end{lemma}

\begin{proof}
映射
\[
\Op^{\mer}_{\check G}\times_{\LS^{\mer}_{\check G}}\check{\mathfrak{n}}/\check B
\longrightarrow \Op^{\mer}_{\check G}
\]
是
\[
\widetilde{\check N}/\check G\longrightarrow \check N/\check G
\longrightarrow \LS^{\mer}_{\check G}
\]
的基变换。第二个映射是闭嵌入，故能检测相干对象。因此只需对
\[
\widetilde{\check N}/\check G\longrightarrow \check N/\check G
\]
证明该断言；再由 fppf 下降，可归约到对 Springer 消解
\[
\widetilde{\check N}\longrightarrow \check N
\]
证明该断言，后者分解为
\[
\widetilde{\check N}\longrightarrow \check N\times \check G/\check B\longrightarrow \check N.
\]
由于第一个映射是闭嵌入，它能检测相干对象，故只需对
\[
\check G/\check B\longrightarrow \mathrm{pt}
\]
证明该断言。但在 $\check G/\check B$ 上，「$\lambda$-扭整体截面函子共同检测相干对象」这一断言
可由以下事实推出：$\QCoh(\check G/\check B)=\IndCoh(\check G/\check B)$ 由有限个
$\mathcal{O}(\lambda)$ 生成，从而
\[
\QCoh(\check G/\check B)\simeq R\text{-}\mathrm{mod}
\]
\srcpage{110}
其中 $R$ 是光滑真 dg 代数，而遗忘函子
$R\text{-}\mathrm{mod}\to \Vect$ 对应有限个扭整体截面函子之和。但后者显然检测相干对象。
\end{proof}

\begin{proposition}\label{prop:3.4.2.4}
函子 IFLE 保持紧对象。
\end{proposition}

\begin{proof}
$\IKL(G)_{\crit}$ 由 Verma 模 $M^{\check\lambda}$ 紧生成，故只需说明
$\IFLE(M^{\check\lambda})$ 是紧的。由上述引理，只需说明
\[
\bigl(\Op^{\mer}_{\check G}\times_{\LS^{\mer}_{\check G}}
\check{\mathfrak{n}}/\check B\to \Op^{\mer}_{\check G}\bigr)^{*}
\bigl(\IFLE(M^{\check\lambda})(\lambda)\bigr)
\]
是紧的。然而 IFLE 与 B 相容，故在谱侧用 $\lambda$ 扭对应于几何侧与 Wakimoto 模
$J_{\lambda}$ 卷积，即
\[
\IFLE(M^{\check\lambda})(\lambda)=\IFLE(M^{\check\lambda}\star J_{\lambda}).
\]
此外，
\[
\bigl(\Op^{\mer}_{\check G}\times_{\LS^{\mer}_{\check G}}
\check{\mathfrak{n}}/\check B\to \Op^{\mer}_{\check G}\bigr)^{*}\circ \IFLE\colon
\IKL(G)_{\crit}\longrightarrow \IndCoh^{*}(\Op^{\mer}_{\check G})
\]
对应于函子
\[
C^{I}\longrightarrow C\longrightarrow \Whit^{*}(C)
\]
（取 $C=\hat{\mathfrak{g}}\text{-}\mathrm{mod}_{\mathrm{crit}}$），它保持紧对象。
$M^{\check\lambda}$ 是紧的，而与 Wakimoto 模卷积是等价，故 $M^{\check\lambda}\star J_{\lambda}$
也是紧的。证毕。
\end{proof}
